\subsection{与一个自同态相关联的模}

设 A 为一个交换环，设 M 为一个 A-模，并设 u 为 M 的一个 A-自同态。回顾（III，第 538 页），映射 \((p(\mathbf{X}), x) \mapsto p(u) \cdot x\)（从 \(A[X] \times M\) 到 M 中）使 M 成为 A{[}X{]}-模，记作 \(M_{u}\)。再回忆（III，第 538 与 539 页）：若 M{[}X{]} 表示 M 经由从 A 到 A{[}X{]}的标量扩张得到的 A{[}X{]}-模，且若 \(\bar{u}\) 表示由 u 诱导的 M{[}X{]} 的 A{[}X{]}-自同态，则存在 A{[}X{]}-模的正合序列 \(^{1}\)

\[
0 \to \mathbf {M} [ \mathbf {X} ] \stackrel {\psi} {\to} \mathbf {M} [ \mathbf {X} ] \stackrel {\varphi} {\to} \mathbf {M} _ {u} \to 0,\tag{1}
\]

，其中 \(\varphi(p(X) \otimes x) = p(u).x\)，且 \(\psi = X - \bar{u}\)。

一个自同态 \(u\)（A-模 M 的）与一个自同态 \(u'\)（A-模 \(M'\) 的）称为相似的，如果存在一个同构 \(g\)（从 M 到 \(M'\) 上的）使得 \(u' \circ g = g \circ u\)，即（III，第 540 页，命题 19）一个同构 \(g\) 从 \(M_u\) 映到 \(M_{u'}'\) 上。如果 M（相应地，\(M'\)）在有限基 B（相应地，\(B'\)）上是自由的，并且如果 \(M(u)\)（相应地 \(M(u')\)）是 \(u\)（相应地 \(u'\)）关于 B（相应地，\(B'\)）的矩阵，则 \(u\) 与 \(u'\) 相似当且仅当 \(M(u)\) 与 \(M(u')\) 是相似矩阵（II，第 356 页，定义 6）。有限生成自由模的两个相似自同态的特征多项式（III，第 541 页，定义 3）相等（III，第 540 页，命题 19）。

设 K 是一个交换域；则由 K 上向量空间 E 与 E 的自同态 u 组成的每个偶对 \((E, u)\) 都对应着一个 K{[}X{]}-模 \(E_{u}\)。由于环 K{[}X{]} 是主理想整环（IV，第 11 页，命题 11），前面各节的结果可以应用于 \(E_{u}\)。

让我们首先说明如何将某些概念从模的语言翻译为向量空间自同态的语言：

« V 是 \(E_{u}\) 的一个子模 » 表示：「V 是 E 的一个在 u 下封闭的向量子空间」。

« V 是 \(E_{u}\) 的一个循环子模 » 意指：「存在 \(x \in V\)，使得向量子空间 V 由元素 \(u^{i} \cdot x (i \in N)\) 张成」。此时称 V 是（关于 u）循环的，并称 x 为一个生成元。

« V 是 \(E_{u}\) 的一个不可分解子模 » 表示：「V 非零，且 V 不是两个各自在 u 下封闭的非零子空间的直和」。

« a 是子模 V 的零化子 » 意指：「a 是由那些多项式 \(p(\mathbf{X}) \in \mathbf{K}[\mathbf{X}]\)（它们满足 \(p(u).x = 0\)，对所有 \(x \in V\)）构成的理想 »。

使 a 等于主理想 \((g)\) 的首一多项式 g 则称为 u 在 V 上限制的最小多项式。

« \(E_{u}\) 是循环的，且零化子为 \(a = (g)\) »

\[
\left(\text { with } g (\mathbf {X}) = \mathbf {X} ^ {n} + \alpha_ {n - 1} \mathbf {X} ^ {n - 1} + \dots + \alpha_ {0}\right)
\]

\(^{1}\) \(\psi\) 的单射性在 III，第 539 页，命题 18 中并未陈述，其证明如下：在引文所用记号下，我们有

\[
\psi \left(\sum \left(\mathrm{X} ^ {k} \otimes x _ {k}\right)\right) = \sum \mathrm{X} ^ {k} \otimes \left(x _ {k - 1} - u (x _ {k})\right).
\]

如果 \(\sum X^k \otimes x_k\) 属于 \(\psi\) 的核，则由此可得 \(x_{k-1} = u(x_k)\) 对所有 \(k\)，从而诸 \(x_k\) 全为零，因为族 \((x_k)\) 具有有限支撑。

表示：「存在 \(x \in E\)，使得 \((u^{i} \cdot x)\) ( \(0 \leqslant i \leqslant n - 1\) ) 是向量空间 E 的一组基，且 \(g(u) \cdot x = 0\)」。换言之，可以找到 E 的一组基，使 u 关于这组基的矩阵 U 为

\[
U = \left( \begin{array}{c c c c} 0 & 0 & 0 \dots 0 & - \alpha_ {0} \\ 1 & 0 & 0 \dots 0 & - \alpha_ {1} \\ 0 & 1 & 0 \dots 0 & - \alpha_ {2} \\ \cdots & \cdots & \cdots & \cdots \\ 0 & 0 & 0 \dots 0 & - \alpha_ {n - 2} \\ 0 & 0 & 0 \dots 1 & - \alpha_ {n - 1} \end{array} \right).\tag{2}
\]

« \(E_{u}\) 是挠模 » 意味着，依照上文给出的循环挠模的刻画：「\(E_{u}\) 的每个循环子模在 K 上都是有限维的」。特别地：

« \(E_{u}\) 是有限生成的挠模 » 意味着：「E 在 K 上是有限维的」。

